\documentclass[11pt]{article}
\usepackage[margin=1in]{geometry}
\usepackage{amsmath,amssymb,amsthm,booktabs,hyperref,microtype}
\hypersetup{colorlinks=true,linkcolor=blue,pdftitle={Uniform Vacuum Weights and Harnack Control},pdfauthor={Yang--Mills representational dynamics research programme}}
\newtheorem{theorem}{Theorem}
\newtheorem{corollary}{Corollary}
\newtheorem{warning}{Warning}
\title{Uniform Vacuum Weights and Harnack Control\\Target 55 of the H1 Downstream Programme}
\author{Yang--Mills representational dynamics research programme}
\date{21 July 2026}
\begin{document}
\maketitle
\begin{abstract}
Finite Perron positivity does not control the inverse vacuum weight uniformly.
We identify a sharp elementary sufficient hypothesis: an entrywise positive
transfer amplitude with uniform lower and upper bounds.  It yields an explicit
Harnack ratio and root-moment bound.  A positive two-sector family then shows
exponential loss when the lower bound degenerates.  This isolates the physical
Yang--Mills obligation as a uniform comparison principle rather than a finite
eigenvector computation.
\end{abstract}

\section{Finite Harnack theorem}
Let $T$ be an $N\times N$ strictly positive matrix and let
$T\psi=\lambda_0\psi$ with $\psi>0$.
\begin{theorem}[Entrywise Harnack bound]
If $0<m\leq T_{ij}\leq M<\infty$ for every $i,j$, then
\begin{equation}\label{eq:harnack}
 \frac{\max_i\psi_i}{\min_i\psi_i}\leq\frac{M}{m}.
\end{equation}
\end{theorem}
\begin{proof}
For any rows $i,k$,
\[
 \frac{\psi_i}{\psi_k}
 =\frac{\sum_jT_{ij}\psi_j}{\sum_jT_{kj}\psi_j}
 \leq\frac{M\sum_j\psi_j}{m\sum_j\psi_j}=\frac Mm.
\]
Interchange $i,k$ and maximize.
\end{proof}

Normalize the forward cost by $V_i=(\max_j\psi_j)/\psi_i$.
\begin{corollary}[Uniform root moment]
$1\leq V_i\leq M/m$, and hence $\nu(V)\leq M/m$ for every probability root
law $\nu$.
\end{corollary}
The associated Doob kernel
\[
 P_{ij}=\frac{T_{ij}\psi_j}{\lambda_0\psi_i}
\]
is stochastic.  Thus the same hypothesis controls both the normalization of
the dynamics and the forward inverse-amplitude cost.

\section{Finite calibration and hostile family}
For the displayed three-state positive matrix in the executable certificate,
the actual amplitude ratio is $1.1222312626$, while the entrywise theorem gives
the conservative bound $5$.

Now set $\epsilon_L=e^{-L}$ and
\[
 T_L=\begin{pmatrix}1&\epsilon_L\\
 \epsilon_L&2\epsilon_L^2\end{pmatrix}.
\]
Every $T_L$ is strictly positive and positive definite for sufficiently large
$L$, yet its suppressed-sector inverse weight grows like $e^L$.  At $L=20$
the actual ratio and hostile root moment are approximately
$4.851651954\times10^8$, while the entrywise bound has deteriorated to
$1.176926334\times10^{17}$.

\begin{warning}
Finite-volume Krein--Rutman or Perron--Frobenius positivity cannot be promoted
to a cutoff-uniform root moment.  A global entrywise minorization may itself be
too strong for Yang--Mills; a chart-chain, sector-aware Harnack inequality or a
direct weighted moment estimate would suffice, but must be proved uniformly.
\end{warning}

\section{Partition-weight diagnostic}
Given boundary partition weights $Z_i$, the normalized vector $Z_i^{-1}$ can
be compared with $\psi_i^{-1}$ by the oscillation of their pointwise ratio.
This is a diagnostic only: equality requires a theorem identifying the
temporal transfer amplitude with the retained conditional disintegration.

\section{Audit}
The Harnack and root-moment implications are C2 finite theorems.  No
cutoff-uniform physical comparison, basin entry, continuum construction, or
gap follows.  Ten hostile tests verify the eigenvector, Doob, Harnack,
root-moment, mismatch, and nonuniformity claims.
\end{document}
